\section{Treatment Knowledge Base}
\label{sec:kb}

A prediction alone does not tell a farmer what to do next, so each of the 38 PlantVillage classes is linked to a short, sourced treatment record in \texttt{knowledge\_base/treatments.json}. The app will look up the predicted class and display that record.

\subsection{Schema}
\texttt{treatments.json} is a single JSON object. Its keys are the exact class strings from \texttt{splits/class\_names.json}, plus one top-level \texttt{disclaimer} string. Every entry is validated against \texttt{knowledge\_base/schema.json} (JSON Schema draft-07), which requires the following 13 fields and rejects any others:
\begin{itemize}
    \item \texttt{crop}, \texttt{disease}, \texttt{is\_healthy} (boolean), \texttt{cause}, \texttt{symptoms};
    \item \texttt{chemical\_treatment}, \texttt{organic\_treatment}, \texttt{immediate\_action}, \texttt{prevention}, \texttt{severity\_guidance};
    \item \texttt{sources} (non-empty list of unique citation strings), \texttt{last\_reviewed} (date), \texttt{reviewed\_by}.
\end{itemize}
Healthy classes use the same fields, with \texttt{chemical\_treatment} set to ``None needed.'' and routine scouting and prevention advice instead of treatment.

\subsection{Validation}
\texttt{data\_scripts/validate\_treatments.py} checks that:
\begin{itemize}
    \item every entry passes the schema, and every text field and source string is non-empty;
    \item every key is a valid class name (no unknown or misspelled classes);
    \item the top-level disclaimer is present;
    \item healthy classes do not recommend a chemical treatment (reported as a warning).
\end{itemize}
It also prints a per-crop summary and the list of classes still missing. Run with \texttt{--require-complete}, it fails unless all 38 classes are present; this is the check for final submission.

\subsection{Example Entry}
Table~\ref{tab:kb_example} shows the entry for \texttt{Tomato\_\_\_Early\_blight}.

\begin{table}[htbp]
\centering
\small
\caption{Knowledge base entry for \texttt{Tomato\_\_\_Early\_blight}.}
\label{tab:kb_example}
\begin{tabular}{>{\raggedright\arraybackslash}p{0.22\linewidth}>{\raggedright\arraybackslash}p{0.70\linewidth}}
\toprule
\textbf{Field} & \textbf{Value} \\
\midrule
\texttt{cause} & Fungus \textit{Alternaria solani} (now also classified as \textit{A.\ linariae}) \\
\texttt{symptoms} & Dark concentric-ringed spots on older leaves, yellowing around spots; can progress up the canopy and affect stems/fruit. \\
\texttt{chemical\_treatment} & Protectant fungicides such as chlorothalonil or mancozeb (azoxystrobin is another option), applied every 7--10 days to protect new growth; do not rely on strobilurins alone because resistance has been reported (follow the product label). \\
\texttt{organic\_treatment} & Fixed copper (copper hydroxide or copper sulfate) or \textit{Bacillus subtilis} products labeled for organic use. \\
\texttt{immediate\_action} & Remove infected plant material from the field; avoid overhead watering. \\
\texttt{prevention} & Rotate away from tomato for 2--3 years; use pathogen-free seed/transplants and resistant or tolerant varieties; stake or trellis for airflow; mulch and prune the lowest leaves to reduce soil splash. \\
\texttt{severity\_guidance} & When lesions merge and lower leaves drop, disease is severe: keep protectant sprays on schedule within label limits and remove infected debris. \\
\texttt{sources} & Connecticut Agricultural Experiment Station, \textit{Early Blight of Tomato} (2012); NC State Extension, \textit{Early Blight of Tomato}. \\
\texttt{last\_reviewed} & 2026-10-12 (Babar) \\
\bottomrule
\end{tabular}
\end{table}

\subsection{Sources}
Every entry cites at least one university extension or plant-protection fact sheet, and all citations are listed with URLs and access dates in \texttt{knowledge\_base/SOURCES.md}. Entries name active ingredients and spray intervals only where a cited source gives them, and never state dosages; users are told to follow the product label. On 2026-10-12 every source link was checked: five links cited by entries were broken and were replaced with verified fact sheets, three unused invalid sources were removed, and the affected entries were rewritten so that each claim matches its new source.

\subsection{Status and Limitations}
\begin{itemize}
    \item \textbf{Coverage:} All 38 entries are 100\% complete, verified against \texttt{knowledge\_base/schema.json}, and validated with zero errors by \texttt{data\_scripts/validate\_treatments.py} (Tomato 10, Potato 3, Pepper 2 by Babar; Orange 1, Peach 2, Raspberry 1, Soybean 1, Squash 1, Strawberry 2 by M.~Anas; Apple 4, Corn 4, Grape 4, Blueberry 1, Cherry 2 by Abdul Ahad).
    \item \textbf{Not region-specific:} The primary sources are extension services (e.g., UC Davis, NC State, Cornell, Florida Extension). Active chemical ingredients and organic remedies are cross-referenced with standard international plant protection guidelines; local agronomist consultation remains recommended.
    \item \textbf{Disclaimer:} The file carries the mandatory top-level disclaimer ``Decision support only. Consult a local agronomist before applying any product. Always follow product labels and local regulations.'', which the app renders alongside every recommendation.
\end{itemize}
